\chapter{Indicadores estadísticos del sistema educativo y demanda laboral (SEP e IMCO 2026)}\label{anexo:estadisticas}

El presente anexo consolida la evidencia estadística nacional que fundamenta la necesidad de optimizar los mecanismos de orientación vocacional en el nivel superior, articulando los datos del Sistema Educativo Nacional y los estudios de mercado laboral elaborados por el Instituto Mexicano para la Competitividad (IMCO).

\section{Cobertura y deserción en la educación superior en México}
De acuerdo con el informe oficial de la Secretaría de Educación Pública~\cite{SEP2026Cifras}, durante el ciclo escolar reciente la cobertura de educación superior en México alcanzó el 44.5\%. Sin embargo, las investigaciones sobre trayectorias estudiantiles revelan que el acceso ampliado no se traduce de forma directa en eficiencia terminal. Diversos estudios sistemáticos~\cite{Villegas2024,KattMorales2024} documentan que entre el 25\% y el 35\% de los estudiantes universitarios abandonan o interrumpen sus estudios durante los primeros tres semestres de la carrera.

Entre los factores determinantes de este fenómeno, las investigaciones pedagógicas identifican tres causas prioritarias:
\begin{enumerate}
    \item \textbf{Incompatibilidad vocacional temprana:} Discrepancia entre las expectativas preconcebidas del estudiante y el rigor técnico real del plan de estudios.
    \item \textbf{Carencia de orientación personalizada:} Ausencia de herramientas cuantitativas que contrasten las habilidades previas del alumno contra los perfiles de ingreso formales de la institución.
    \item \textbf{Desinformación respecto a la oferta formativa:} Concentración de la demanda en carreras de alta visibilidad mediática, ignorando programas técnicos especializados con mayor afinidad individual.
\end{enumerate}

\section{Concentración de matrícula universitaria según el IMCO}
El análisis sectorial desarrollado por el Instituto Mexicano para la Competitividad en su plataforma \textit{Compara Carreras 2026}~\cite{IMCO2026} revela un desbalance estructural crítico: a nivel nacional, el 46\% de la matrícula universitaria total se concentra en apenas diez programas académicos tradicionales. La tabla~\ref{tab:imco-concentracion} sintetiza la distribución de estos diez programas, evidenciando las altas tasas de saturación laboral e informalidad que enfrentan sus egresados en comparación con áreas de ciencia y tecnología.

\begin{table}[htbp]
\centering
\small
\renewcommand{\arraystretch}{1.2}
\caption{Concentración de matrícula y condiciones laborales en las diez carreras más demandadas de México~\cite{IMCO2026}}
\label{tab:imco-concentracion}
\begin{tabularx}{\textwidth}{L{0.28\textwidth}L{0.18\textwidth}L{0.24\textwidth}X}
\toprule
\rowcolor{encabezadoTabla} \color{white}\textbf{Programa académico} & \color{white}\textbf{Participación} & \color{white}\textbf{Tasa de informalidad} & \color{white}\textbf{Riesgo de saturación} \\
\midrule
1. Administración de empresas & 8.6\% & 24.1\% & Muy alto \\
\rowcolor{filaAlterna}
2. Derecho & 8.2\% & 27.5\% & Muy alto \\
3. Contabilidad y fiscalización & 6.1\% & 21.8\% & Alto \\
\rowcolor{filaAlterna}
4. Psicología & 4.9\% & 32.4\% & Muy alto \\
5. Ingeniería industrial & 4.5\% & 18.2\% & Medio \\
\rowcolor{filaAlterna}
6. Medicina general & 4.1\% & 14.3\% & Competencia extrema \\
7. Enfermería general & 3.4\% & 19.6\% & Medio a alto \\
\rowcolor{filaAlterna}
8. Ciencias de la comunicación & 2.4\% & 36.8\% & Crítico \\
9. Arquitectura y urbanismo & 2.1\% & 29.4\% & Alto \\
\rowcolor{filaAlterna}
10. Ingeniería mecánica & 1.7\% & 16.5\% & Medio \\
\midrule
\textbf{Total acumulado (Top 10)} & \textbf{46.0\%} & \textbf{Promedio: 24.1\%} & \textbf{Concentración estructural} \\
\bottomrule
\end{tabularx}
\end{table}

\section{Oportunidad estratégica en la oferta académica del IPN}
En contraste con este patrón de saturación nacional, el Instituto Politécnico Nacional en planteles de la Ciudad de México ofrece programas académicos altamente especializados en las ramas de Ingeniería y Ciencias Físico-Matemáticas, Ciencias Médico-Biológicas y Ciencias Sociales y Administrativas (por ejemplo, Ingeniería Telemática, Ingeniería Biónica, Ingeniería en Sistemas Automotrices, Inteligencia Artificial y Farmacia). 

No obstante, la falta de herramientas que evalúen competencias analíticas y técnicas en el aspirante conduce a que muchos de ellos elijan únicamente por denominación popular. La recomendación vocacional basada en aprendizaje automático y evidencia textual adquiere relevancia nacional al visibilizar estas alternativas académicas pertinentes, guiando al aspirante hacia opciones formativas con correspondencia matemática verificable y menor riesgo de deserción escolar.

\endinput
